% Options for packages loaded elsewhere
\PassOptionsToPackage{unicode}{hyperref}
\PassOptionsToPackage{hyphens}{url}
%
\documentclass[
]{article}
\usepackage{amsmath,amssymb}
\usepackage{iftex}
\ifPDFTeX
  \usepackage[T1]{fontenc}
  \usepackage[utf8]{inputenc}
  \usepackage{textcomp} % provide euro and other symbols
\else % if luatex or xetex
  \usepackage{unicode-math} % this also loads fontspec
  \defaultfontfeatures{Scale=MatchLowercase}
  \defaultfontfeatures[\rmfamily]{Ligatures=TeX,Scale=1}
\fi
\usepackage{lmodern}
\ifPDFTeX\else
  % xetex/luatex font selection
\fi
% Use upquote if available, for straight quotes in verbatim environments
\IfFileExists{upquote.sty}{\usepackage{upquote}}{}
\IfFileExists{microtype.sty}{% use microtype if available
  \usepackage[]{microtype}
  \UseMicrotypeSet[protrusion]{basicmath} % disable protrusion for tt fonts
}{}
\makeatletter
\@ifundefined{KOMAClassName}{% if non-KOMA class
  \IfFileExists{parskip.sty}{%
    \usepackage{parskip}
  }{% else
    \setlength{\parindent}{0pt}
    \setlength{\parskip}{6pt plus 2pt minus 1pt}}
}{% if KOMA class
  \KOMAoptions{parskip=half}}
\makeatother
\usepackage{xcolor}
\usepackage{color}
\usepackage{fancyvrb}
\newcommand{\VerbBar}{|}
\newcommand{\VERB}{\Verb[commandchars=\\\{\}]}
\DefineVerbatimEnvironment{Highlighting}{Verbatim}{commandchars=\\\{\}}
% Add ',fontsize=\small' for more characters per line
\newenvironment{Shaded}{}{}
\newcommand{\AlertTok}[1]{\textcolor[rgb]{1.00,0.00,0.00}{\textbf{#1}}}
\newcommand{\AnnotationTok}[1]{\textcolor[rgb]{0.38,0.63,0.69}{\textbf{\textit{#1}}}}
\newcommand{\AttributeTok}[1]{\textcolor[rgb]{0.49,0.56,0.16}{#1}}
\newcommand{\BaseNTok}[1]{\textcolor[rgb]{0.25,0.63,0.44}{#1}}
\newcommand{\BuiltInTok}[1]{\textcolor[rgb]{0.00,0.50,0.00}{#1}}
\newcommand{\CharTok}[1]{\textcolor[rgb]{0.25,0.44,0.63}{#1}}
\newcommand{\CommentTok}[1]{\textcolor[rgb]{0.38,0.63,0.69}{\textit{#1}}}
\newcommand{\CommentVarTok}[1]{\textcolor[rgb]{0.38,0.63,0.69}{\textbf{\textit{#1}}}}
\newcommand{\ConstantTok}[1]{\textcolor[rgb]{0.53,0.00,0.00}{#1}}
\newcommand{\ControlFlowTok}[1]{\textcolor[rgb]{0.00,0.44,0.13}{\textbf{#1}}}
\newcommand{\DataTypeTok}[1]{\textcolor[rgb]{0.56,0.13,0.00}{#1}}
\newcommand{\DecValTok}[1]{\textcolor[rgb]{0.25,0.63,0.44}{#1}}
\newcommand{\DocumentationTok}[1]{\textcolor[rgb]{0.73,0.13,0.13}{\textit{#1}}}
\newcommand{\ErrorTok}[1]{\textcolor[rgb]{1.00,0.00,0.00}{\textbf{#1}}}
\newcommand{\ExtensionTok}[1]{#1}
\newcommand{\FloatTok}[1]{\textcolor[rgb]{0.25,0.63,0.44}{#1}}
\newcommand{\FunctionTok}[1]{\textcolor[rgb]{0.02,0.16,0.49}{#1}}
\newcommand{\ImportTok}[1]{\textcolor[rgb]{0.00,0.50,0.00}{\textbf{#1}}}
\newcommand{\InformationTok}[1]{\textcolor[rgb]{0.38,0.63,0.69}{\textbf{\textit{#1}}}}
\newcommand{\KeywordTok}[1]{\textcolor[rgb]{0.00,0.44,0.13}{\textbf{#1}}}
\newcommand{\NormalTok}[1]{#1}
\newcommand{\OperatorTok}[1]{\textcolor[rgb]{0.40,0.40,0.40}{#1}}
\newcommand{\OtherTok}[1]{\textcolor[rgb]{0.00,0.44,0.13}{#1}}
\newcommand{\PreprocessorTok}[1]{\textcolor[rgb]{0.74,0.48,0.00}{#1}}
\newcommand{\RegionMarkerTok}[1]{#1}
\newcommand{\SpecialCharTok}[1]{\textcolor[rgb]{0.25,0.44,0.63}{#1}}
\newcommand{\SpecialStringTok}[1]{\textcolor[rgb]{0.73,0.40,0.53}{#1}}
\newcommand{\StringTok}[1]{\textcolor[rgb]{0.25,0.44,0.63}{#1}}
\newcommand{\VariableTok}[1]{\textcolor[rgb]{0.10,0.09,0.49}{#1}}
\newcommand{\VerbatimStringTok}[1]{\textcolor[rgb]{0.25,0.44,0.63}{#1}}
\newcommand{\WarningTok}[1]{\textcolor[rgb]{0.38,0.63,0.69}{\textbf{\textit{#1}}}}
\usepackage{graphicx}
\makeatletter
\def\maxwidth{\ifdim\Gin@nat@width>\linewidth\linewidth\else\Gin@nat@width\fi}
\def\maxheight{\ifdim\Gin@nat@height>\textheight\textheight\else\Gin@nat@height\fi}
\makeatother
% Scale images if necessary, so that they will not overflow the page
% margins by default, and it is still possible to overwrite the defaults
% using explicit options in \includegraphics[width, height, ...]{}
\setkeys{Gin}{width=\maxwidth,height=\maxheight,keepaspectratio}
% Set default figure placement to htbp
\makeatletter
\def\fps@figure{htbp}
\makeatother
\setlength{\emergencystretch}{3em} % prevent overfull lines
\providecommand{\tightlist}{%
  \setlength{\itemsep}{0pt}\setlength{\parskip}{0pt}}
\setcounter{secnumdepth}{-\maxdimen} % remove section numbering
\ifLuaTeX
\usepackage[bidi=basic]{babel}
\else
\usepackage[bidi=default]{babel}
\fi
\babelprovide[main,import]{chinese}
% get rid of language-specific shorthands (see #6817):
\let\LanguageShortHands\languageshorthands
\def\languageshorthands#1{}
\ifLuaTeX
  \usepackage{selnolig}  % disable illegal ligatures
\fi
\usepackage{bookmark}
\IfFileExists{xurl.sty}{\usepackage{xurl}}{} % add URL line breaks if available
\urlstyle{same}
\hypersetup{
  pdftitle={AB 云：相位谐振腔},
  pdfauthor={Isaev Iskhak Khamzatovich（ORCID: 0009-0003-7299-0701）},
  pdflang={zh-CN},
  pdfkeywords={希尔伯特--波利亚猜想, 黎曼 ζ
函数, 随机矩阵理论, GUE, 阿哈罗诺夫--玻姆效应, Hofstadter 哈密顿量, AB
云, 拓扑涡旋, 狄拉克锥, TKNN},
  hidelinks,
  pdfcreator={LaTeX via pandoc}}

\title{AB 云：相位谐振腔}
\usepackage{etoolbox}
\makeatletter
\providecommand{\subtitle}[1]{% add subtitle to \maketitle
  \apptocmd{\@title}{\par {\large #1 \par}}{}{}
}
\makeatother
\subtitle{黎曼 ζ 零点、GUE
普适性与拓扑涡旋物质------完整数值研究与独立验证}
\author{Isaev Iskhak Khamzatovich（ORCID: 0009-0003-7299-0701）}
\date{第 22 版 · 2026 年 8 月}

\begin{document}
\maketitle

\section{摘要}\label{ux6458ux8981}

本专著从零重新撰写，系统总结了我们称之为 \textbf{AB
云}的数值研究：这是一种格点量子系统（带阿哈罗诺夫--玻姆相位的 Hofstadter
哈密顿量），其中拓扑涡旋激发产生非均匀相位织构，其本征值分布与黎曼 ζ
函数零点的统计进行比较。文中每一个论断都附有自动化的 37
项测试验证套件（v18/v19 版，2026 年 8
月）给出的可复现数字：每个数字都标注测试编号，每项测试都生成完整计算日志，每次套件运行都会在磁盘上留下包含四种格式报告（Markdown、HTML、PDF、DOCX）、三种图形格式以及逐时间戳日志的目录，供独立验证之用。

主要验证结果如下。第一，在光滑规范（\texttt{:monumental}
涡旋相位模型）下，环面上 AB 云谱的相邻能隙比平均值为
\(\langle r\rangle = 0.585 \pm 0.026\)，GUE 目标值为 \(0.5992\)，偏差
\(-2.4\%\)（测试 33）；对晶格尺寸 \(L = 10, 20, 30, 50\)
的扫描显示单调收敛 \(0.548 \to 0.592 \to 0.605 \to 0.595\)（测试
34）。第二，将 AB 云的展析间距与最初 5000 个 ζ
零点直接比较，云的对关联函数在整个区间 \(s\in[0;3]\) 上更接近 GUE
而非泊松，但在如此小的系综下，双样本 Kolmogorov--Smirnov
检验拒绝分布完全相同的假设（\(p = 6.7\cdot10^{-4}\)；测试
35）------我们如实记录该偏差并以系综有限性加以解释。第三，构造的拓扑骨架以机器精度得到确认：狄拉克弦规范下恰好有一个元胞携带磁通
\(2\pi q\)（测试 15、24）；Byers--Yang 定理对整数电荷成立到
\(3.5\cdot10^{-15}\)，分数电荷则偏离 \(0.34\)（测试
25）；最低朗道能级的陈数为 \(C_1 = 2\)（测试 20）；\(\alpha = 1/2\) 处的
Connes 自对偶产生恰好 4 个零模（测试
17）。最后，低能动力学呈现狄拉克锥，能隙线性标度
\(E_{\min}\propto 1/L\)，\(R^2 = 0.9997\)（测试 19）；\(\alpha = 1/2\)
处态密度呈现对比度 \(20\times\) 的狄拉克凹陷（测试 30）；偏离厄米点
\(\sigma = 1/2\) 时出现 Hatano--Nelson 趋肤效应（测试 32）。

这些结果共同构成连贯的图景：AB 云是一个相位谐振腔，普适的 GUE
统计源于复数跃迁相位对时间反演对称性的破坏，涡旋表现得像相对论性粒子，而临界线
\(\sigma = 1/2\) 作为自对偶点格外突出------云的统计在此处最接近 ζ
零点的统计。我们始终区分''已被证明的（机器精度）``、''可靠测量的（校准统计带）``与''仍属开放的（大
\(T\)
渐近）``，并为每个结论提供明确判据------读者无需盲信本书中的任何一个数字。

\textbf{关键词：} 希尔伯特--波利亚猜想；黎曼 ζ
函数；随机矩阵理论；GUE；阿哈罗诺夫--玻姆效应；Hofstadter 哈密顿量；AB
云；拓扑涡旋；狄拉克锥；TKNN；Hatano--Nelson 趋肤效应。

\section{1.
引言与问题设定}\label{ux5f15ux8a00ux4e0eux95eeux9898ux8bbeux5b9a}

\subsection{1.1
希尔伯特--波利亚猜想}\label{ux5e0cux5c14ux4f2fux7279ux6ce2ux5229ux4e9aux731cux60f3}

1914 年黎曼猜想 ζ 函数的所有非平凡零点位于临界线
\(\mathrm{Re}\,s = 1/2\) 上。近半个世纪后的 1958 年，希尔伯特（据 Pólya
回忆）提出了如今称为\textbf{希尔伯特--波利亚猜想}的策略：若能找到谱与零点纵坐标
\(\{\gamma_n\}\) 重合的自共轭算符
\(\hat H\)，则黎曼假设立即成为定理------厄米性迫使所有本征值为实，即对所有
\(n\) 有
\(\mathrm{Re}\,\gamma_n = 1/2\)。该纲领的美妙之处在于把解析数论翻译为量子力学语言：ζ
零点成为某个（尚未知晓的）物理系统的能级。

该猜想至今未被证明，但它恰好构成本书的框架：我们构造并研究一类具体的量子格点系统，其谱统计逼近普适的随机矩阵行为。必须强调区别于以往文献的方法论立场：我们\textbf{不宣称证明黎曼假设}。我们的对象是对''有限格点系综如何逼近普适随机矩阵统计、哪些参数使逼近更近或更远''的精确数值刻画。所有''逼近程度''的判据在运行测试之前预先固定，研究中绝不调整。

历史背景简述：1972 年 Montgomery 提出零点对关联猜想；Dyson
在普林斯顿高等研究院茶室偶遇 Montgomery 时指出，该公式与随机矩阵理论中
GUE（高斯幺正系综）本征值对关联一致。1980 年代 Odlyzko
在庞大数据上数值确认了这一吻合；1985 年 Berry
计算了有限尺寸修正：真实的（有限 \(T\)）零点样本必须以可预测方式偏离纯
GUE 预测------特别是原点处对关联 \(R_2(0)\) 获得按 \(T\)
幂次衰减的修正。这些修正对我们至关重要：它们解释了为什么测试 4--6 与 11
在 \(T\lesssim 4\cdot10^4\) 的样本上看到统计显著的 GUE
偏离，以及为什么这种偏离是证实而非否定整体图景。

\subsection{1.2 阿哈罗诺夫--玻姆效应与 AB
云}\label{ux963fux54c8ux7f57ux8bfaux592bux73bbux59c6ux6548ux5e94ux4e0e-ab-ux4e91}

\textbf{阿哈罗诺夫--玻姆效应（AB
效应）}指出：量子粒子即使在场强为零的区域也对磁通有响应：环绕携带磁通
\(\Phi\) 的螺线管的波函数获得相位
\(\varphi_{\rm AB} = 2\pi\Phi/\Phi_0\)，其中 \(\Phi_0 = h/e\)
为磁通量子。在格点上，这通过跃迁振幅的相位修饰实现：\(t_{ij}\to t_{ij}e^{i\varphi_{ij}}\)。\textbf{Hofstadter
哈密顿量}是经典情形：正方格点上的电子处于每元胞磁通
\(\alpha = \Phi/\Phi_0\) 的均匀磁场中；其谱作为 \(\alpha\)
的函数构成著名的 Hofstadter 蝴蝶（见我们的图 19）。

\textbf{AB
云}是该模型的推广：在格点上放置点状\textbf{拓扑涡旋}------相位绕其旋转
\(2\pi q\) 的中心（\(q\)
为涡旋荷，可为分数）。这种系统的相位场是非均匀的：每个涡旋附近相位''卷绕''，整体图像形如云团------名称由此而来。数学上哈密顿量为

\[H_{ij} = -\exp\bigl(i\varphi_{ij}\bigr),\qquad \varphi_{ij} = 2\pi\alpha\,j + \sum_{k=1}^{N_v} q_k\bigl[\arg(\mathbf r_i - \mathbf r_k) - \arg(\mathbf r_j - \mathbf r_k)\bigr] + \varphi^{\rm str}_{ij},\]

其中 \(j\) 为格点列号（Landau 规范），\(N_v\)
为涡旋数，\(\varphi^{\rm str}_{ij}\)
为显式\textbf{狄拉克弦}：穿过涡旋出发的半直线（弦）的每条键附加
\(-2\pi q\)。狄拉克弦保证涡旋周围的总磁通以机器精度等于
\(2\pi q\)------已由直接计算验证（测试 15、24）。

除几何（相位）部分外，模型还包含\textbf{格点无序}
\(W\)：模拟真实晶体杂质的格点加性噪声。模型参数：晶格尺寸
\(L\times L\)、磁通 \(\alpha\)、电荷集合 \(\{q_k\}\)、涡旋数目与位置
\(N_v\)、无序强度
\(W\)、边界条件（开边界或环面）。关键的物理开关是\textbf{涡旋相位模型}：（a）狄拉克弦规范------整数荷涡旋在谱上不可见（这是精确的
Byers--Yang 定理，测试 25）；（b）光滑规范 \texttt{:monumental}（atan
模型）------相位平滑分布且对任意 \(q\) 均为复数------这正是 GUE
统计在物理合理参数下可达的原因。两种规范均在代码中实现并按用途使用：磁通定理在第一种规范中验证，谱统计在第二种规范中测量。

\subsection{1.3
与随机矩阵理论的联系}\label{ux4e0eux968fux673aux77e9ux9635ux7406ux8bbaux7684ux8054ux7cfb}

\textbf{随机矩阵理论（RMT）}研究随机大矩阵的本征值分布。对我们重要的是按时间反演对称性区分的三个系综：\textbf{GOE}（实对称，T
对称）、\textbf{GUE}（复厄米，T
破缺）、\textbf{GSE}（四元数自对偶），以及极限的\textbf{泊松}情形（无关联能级；可积系统）。诊断量定义见
2.2
节；这里只给出相邻能隙比的中心标度：\(\langle r\rangle \approx 0.386\)（泊松）、\(0.531\)（GOE）、\(0.600\)（GUE）。Dyson
三分法预言：复相位对 T 对称性的破坏必然把系统送入 GUE
类------光滑规范下的 AB 云正是教科书式例证：\(\alpha = 1/2\)
时矩阵为复厄米，\(\tau(H) = \|H-H^*\|/\|H\| = 0.225 \neq 0\)（测试
26），统计确实趋向 GUE；而在实数的狄拉克弦规范下 \(\tau = 0\)，谱被锁进
GOE 类（测试 15、16）。

\subsection{1.4
目标、任务与验证的组织}\label{ux76eeux6807ux4efbux52a1ux4e0eux9a8cux8bc1ux7684ux7ec4ux7ec7}

本工作的目标是为 AB 云谱物理建立完整、可独立核查的图景。任务具体化为六块
37 项自动化测试：（1）\(b(N)\) 修正的收敛与评审质疑（测试
1--14）；（2）构造与拓扑定理（15、17、20、24--27）；（3）云的 GUE
分类（16、26、29、33、34）；（4）与 ζ
零点的直接比较（28、35、36）；（5）狄拉克动力学与非厄米性（19、30--32、37）；（6）解析诠释（测试
21--23 与第 8
章）。每项测试按预先登记的判据返回布尔裁决；验证为双通（第一通为主晶格，第二通为加大晶格）；``快速检查''在
\(16\times16\to32\times32\) 晶格上、ζ 样本截断为 5000
个零点，数分钟内跑完全部 37 项测试。

我们明确区分\textbf{两套裁决系统}以免读者混淆。公式型测试（恒等式：磁通、自对偶、Byers--Yang）以机器容差
\(10^{-12}\dots10^{-16}\) 检查，给出二元 PASS/FAIL。统计型测试（GUE
一致性）按模拟确立的置信带校准：例如 \(\langle r\rangle\) 的 PASS
要求落在 \(0.5992\pm0.022\) 带内；这类测试允许第三种裁决
WARN------``偏差已测得但可由有限尺寸解释''。双轨制消除了数值工作的原罪------按结果调判据。

\section{2. 数学准备}\label{ux6570ux5b66ux51c6ux5907}

\subsection{2.1 黎曼 ζ
函数及其零点}\label{ux9eceux66fc-ux3b6-ux51fdux6570ux53caux5176ux96f6ux70b9}

\textbf{黎曼 ζ 函数} \(\zeta(s) = \sum_{n\ge1} n^{-s}\)
解析延拓到全复平面，在临界带 \(0<\mathrm{Re}\,s<1\)
内有非平凡零点。数值上我们使用最初 50,000 个零点（Odlyzko
数据集），纵坐标覆盖 \(T\in[14.13;\,40433.7]\)。高度 \(T\)
处的零点密度由\textbf{黎曼--von Mangoldt 公式}给出：

\[N(T) = \frac{T}{2\pi}\ln\frac{T}{2\pi} - \frac{T}{2\pi} + \frac{7}{8} + O(\ln T),\]

大致按 \((T/2\pi)\ln T\)
增长。除以局部密度------\textbf{展析（unfolding）}------使归一化间距
\(\delta_n = \gamma_{n+1}^{(u)} - \gamma_n^{(u)}\) 均值为
1，全部关联信息随之进入其分布。代码中展析用对光滑函数 \(N(T)\)
的对称滑窗实现；这一标准步骤也是伪影最易混入之处，因此测试 10
检查所有最终量对子样本的稳定性（最大相对偏差 8.5\%，阈值
10\%------PASS）。

我们在三个高度使用零点统计：完整 50,000 零点样本（测试
4、6--10）、\(T>T_{\min}\) 的''高温''子样本（测试
5）、以及用于双通比较的最初 1000/5000 个零点（测试 28、35）。GUE
统计的渐近本性意味着一致性随 \(T\) 改善；我们的数据直接演示了这一点（3.2
节）。

\subsection{2.2
随机矩阵系综与诊断统计量}\label{ux968fux673aux77e9ux9635ux7cfbux7efcux4e0eux8bcaux65adux7edfux8ba1ux91cf}

所有诊断量对展析（无量纲）能级定义。我们使用五个量，各对应关联的不同侧面。

\textbf{相邻能隙比} \(\langle r\rangle\)。对每组相邻三能级
\(E_{n-1},E_n,E_{n+1}\) 令
\(r_n = \min(\delta_n,\delta_{n+1})/\max(\delta_n,\delta_{n+1})\)；其平均无需外部展析（比值对局部密度不敏感）。大矩阵参考值：泊松
\(0.3863\)、GOE \(0.5307\)、GUE
\(0.5996\)。这是有限格点系综的首选诊断量，因为它不积累展析误差。

\textbf{Kolmogorov--Smirnov（KS）检验。} 将间距的经验分布函数与 GUE
Wigner 预测比较：统计量 \(D\) 为两曲线最大绝对差；\(n\) 个间距时 5\%
临界值为 \(1.36/\sqrt{n}\)。\(n = 49999\) 时临界
\(D = 0.0061\)------即使微观偏差也会''拒绝'' GUE；因此我们总是把 \(D\)
与 \(n\)、临界阈值一起报告（测试 4），并单独跟踪趋势
\(D(T_{\min})\)（测试 5）。

\textbf{数目方差 \(\Sigma^2(L)\)}------长为
\(L\)（以平均密度为单位）窗口内能级数的方差。泊松给出
\(\Sigma^2 = L\)；GUE 给出
\(\Sigma^2 \approx (1/\pi^2)\ln(2\pi L) + \ldots\)------对数而非线性增长。谱刚性
\(\Delta_3(L)\) 是其积分版本：窗口 \(L\)
内谱阶梯的最佳直线拟合的均方偏差；泊松 \(L/15\)，GUE
\(\approx (1/\pi^2)\ln L\)。两者均在共同的九点 \(L\) 网格上对数据和模拟
GUE 参考计算（测试 13、14）。

\textbf{对关联 \(R_2(s)\)}------在距给定能级 \(s\)
处找到能级的概率密度（Montgomery 对关联统计量）：GUE 时
\(R_2(s) = 1 - (\sin \pi s/(\pi s))^2\)，原点关联洞 \(R_2(0)=0\)
表达能级排斥；泊松时 \(R_2 \equiv 1\)。谱形因子 \(K(t)\)------\(R_2\)
的傅里叶变换，GUE 为''斜坡+平台''
\(K(t)=\min(t,1)\)------是我们最积分的诊断量（测试 35、36）。

\textbf{Anderson--Darling \(A^2\)}------对分布尾部加权的 KS 类检验；p
值用 2000 次蒙特卡洛模拟估计（测试 11）。最后，\textbf{f\_GUE}
为汇总的''GUE
份额''：\(f_{\rm GUE} = (\Sigma^2_{\rm P} - \Sigma^2_{\rm data})/(\Sigma^2_{\rm P} - \Sigma^2_{\rm GUE})\)，在
\(L=2\) 取值；\(f_{\rm GUE}=1\) 为理想 GUE，\(0\) 为泊松（测试 29）。

\subsection{2.3 带涡旋的 Hofstadter
哈密顿量：精确定义}\label{ux5e26ux6da1ux65cbux7684-hofstadter-ux54c8ux5bc6ux987fux91cfux7cbeux786eux5b9aux4e49}

汇总完整定义。晶格 \(L_x\times L_y\)；波函数定义在格点上；最近邻矩阵元
\(H_{ij} = -\exp(i\varphi_{ij})\)（选择环面时含周期回绕）。Landau
规范给列 \(j\) 相位 \(2\pi\alpha j\)。位于
\(\mathbf r_k=(x_k,y_k)\)（以元胞为单位）的荷 \(q_k\) 涡旋对键相位差
\((i\to j)\) 贡献
\(q_k[\arg(\mathbf r_i-\mathbf r_k)-\arg(\mathbf r_j-\mathbf r_k)]\)，狄拉克弦对穿过弦的键添加离散的
\(-2\pi q_k\)；光滑规范中弦不存在，atan
相位分布于竖直键。关键性质：\textbf{绕任一元胞的相位和等于
\(2\pi\alpha + 2\pi q_k\,\delta_k\)}（再加
\(2\pi\times\)整数），\(\delta_k\) 指示该元胞含涡旋。这已在
\(16\times16\) 晶格的全部 224 个内部元胞上验证：空元胞上
\(|\)磁通\(|\)最大值 \(1.5\cdot10^{-14}\) 弧度（测试 24）。

格点无序 \(W\) 实现为格点势
\(\varepsilon_i\in[-W/2;W/2]\)，生成器种子固定（可复现性：同种子------同无序）。重要的是，对拓扑恒等式（Byers--Yang）无序须关闭：谱对
\(q\) 的任何超出纯规范的依赖都是错误，测试 25 确认到
\(3.5\cdot10^{-15}\)。

边界条件是独立的物理选择。开边界产生蝴蝶的边缘态，混入泊松成分并压低
\(\langle r\rangle\)（测试 33：其余参数相同时开边界 \(0.442\)，环面
\(0.585\)）。环面要求规范相容条件
\(q\,(2-N_x)\in\mathbb{Z}\)------回绕相位周期性；整数 \(q\)
自动满足，\(16\times16\) 环面上的分数 \(q=0.3\)
不满足，系统诚实地回退到开边界（代码自动判别并给出诊断信息）。

\subsection{2.4
验证方法学与可复现性}\label{ux9a8cux8bc1ux65b9ux6cd5ux5b66ux4e0eux53efux590dux73b0ux6027}

套件每次运行创建
\texttt{results/run\_\textless{}时间戳\textgreater{}/}，每项测试一个子目录：内含完整计算日志（每个量带时间戳及相对前项的增量）、控制台输出捕获、Markdown/HTML/矢量
PDF/DOCX 四种报告、PNG（600 dpi）/SVG/PDF 图形，以及首图的逐帧 GIF
动画。日志设计为：外部研究者无需访问我们的代码即可重现本书的每个数字------公式、输入、输出、耗时。复现命令：单测试
\texttt{julia\ ab\_cloud\_v19.jl\ -\/-test\ \textless{}N\textgreater{}\ -\/-no-two-pass}；全套
\texttt{-\/-test\ all}；快速检查
\texttt{-\/-quick}（\(16\times16\to32\times32\)、5000
零点、双通）。报告引擎与计算核心隔离：任一格式失败不能影响测试裁决（v19
架构保证，见附录 A）。

\section{\texorpdfstring{3. 第一块：\(b(N)\) 修正与 ζ 零点统计（测试
1--14）}{3. 第一块：b(N) 修正与 ζ 零点统计（测试 1--14）}}\label{ux7b2cux4e00ux5757bn-ux4feeux6b63ux4e0e-ux3b6-ux96f6ux70b9ux7edfux8ba1ux6d4bux8bd5-114}

\subsection{\texorpdfstring{3.1 什么是
\(b(N)\)，其收敛速率为何重要}{3.1 什么是 b(N)，其收敛速率为何重要}}\label{ux4ec0ux4e48ux662f-bnux5176ux6536ux655bux901fux7387ux4e3aux4f55ux91cdux8981}

修正项 \(b(N)\)
出现于零点对关联的解析处理中：它是把诊断量的有限样本值与其渐近值联系起来的拟合量。程序早期版本曾给它匹配
\(\sim N^{-1/2}\)
的理论预期；评审者正确指出该预期从未被推导。我们接受这一质疑：自套件 v6
起，\(b(N)\)
的指数是\textbf{纯经验}量，没有理论模板；第一块的全部测试检查的是测量的内部一致性，而非与不存在的定理的一致性。测试
3、7、9 提供对同一指数的三种独立测量------它们的一致性正是验证对象。

\subsection{3.2
指数的三种独立测量}\label{ux6307ux6570ux7684ux4e09ux79cdux72ecux7acbux6d4bux91cf}

\begin{figure}
\centering
\includegraphics{../figures/fig01_bN_convergence.png}
\caption{\textbf{图 1.} 修正项 \(b(N)\)
的收敛：数据与两种竞争拟合（幂律与 \(1/\ln N\)）。评审者坚持的
\(N^{-1/2}\) 预期被数据排除。}
\end{figure}

\textbf{对数--对数回归}（测试 3）在网格 \(N\in[50;25600]\) 上给出
\(b(N) \approx 7.03\,N^{-0.1685}\)，\(R^2 = 0.9895\)，斜率标准误
\(0.0043\)。\textbf{延展回归至 \(N=32000\)}（测试 7）给出斜率
\(-0.1504\)，95\% 置信区间
\([-0.1594;-0.1414]\)，\(R^2 = 0.9954\)------相互一致。\textbf{自助法}（测试
9，1000 次有放回重采样）给出中位斜率 \(-0.1746\)，95\% 区间
\([-0.1895;-0.1552]\)。一个通常被忽略的细节：\textbf{替代}参数化
\(b(N) \approx -0.40 + 16.84/\ln N\) 对数据的描述优于幂律（测试 7 中
\(R^2 = 0.9999\) 对
\(0.9954\)）。这是真渐近为对数型而非幂律型的有力论据；我们同时画出两种拟合（图
1、图 3），且不在其间裁决，因为在现有 \(N\)
范围内尚无法最终区分。残差分析（测试 8）显示残差有系列结构（游程数
3，随机期望 5.8；滞后 1 自相关
\(0.68\)）------数据围绕幂律拟合''弯曲''的方式恰如对数形式所规定，与双拟合图像一致。

本块的实践结论：在 \(N\le5\cdot10^4\) 的样本上，\(b(N)\) 保持为 1
的量级（从 \(N=500\) 的 \(2.39\) 到 \(N=32000\) 的 \(1.28\)）------ζ
统计的有限样本效应很大；任何忽略 \(b(N)\) 的''GUE
与否''检验在这些高度上根本无效。

\subsection{\texorpdfstring{3.3 质疑之二：KS、χ² 与 \(A^2\) 拒绝
GUE------而这正是预期的}{3.3 质疑之二：KS、χ² 与 A\^{}2 拒绝 GUE------而这正是预期的}}\label{ux8d28ux7591ux4e4bux4e8cksux3c7uxb2-ux4e0e-a2-ux62d2ux7edd-gueux800cux8fd9ux6b63ux662fux9884ux671fux7684}

\begin{figure}
\centering
\includegraphics{../figures/fig02_spacing_hist.png}
\caption{\textbf{图 2.} 49,999 个 ζ 零点间距的直方图，对 Wigner
猜想曲线（GUE、GOE）与泊松指数：典型的''矩阵型''形状，带能级排斥。}
\end{figure}

\begin{figure}
\centering
\includegraphics{../figures/fig05_ks_convergence.png}
\caption{\textbf{图 3.} KS 统计量 \(D\) 随下截断 \(T_{\min}\)
的变化：向临界线的缓慢漂移------GUE 渐近性的数值表达。}
\end{figure}

测试 4、5、6、11 以最强形式陈述评审质疑并诚实地检验它。在 49,999
个间距的完整样本上，对 GUE 的 KS 统计量 \(D = 0.0881\)，临界值
\(0.0061\)------超出阈值 14 倍，\(H_0\)（GUE）以 \(p\approx10^{-300}\)
被拒。150/300/600 柱的直方图卡方给出 \(\chi^2 = 6170/6276/6384\)（自由度
149/295/550）------所有尺度上均被拒。5000 间距对解析 GUE 分布函数的
Anderson--Darling 给出 \(A^2 = 117.7\)，而蒙特卡洛 95 分位为
\(16.6\)（2000 次模拟 \(p\approx0\)）。

把这读作''对零点 GUE 假设的否定''将是错误的------原因如下。GUE
统计是\textbf{渐近的}：它要在
\(T\gg10^6\)、有数百万零点可用的高度才确立。Berry（1985）计算了 leading
有限尺寸修正，其符号与量级与观测完全吻合：在我们的高度
\(T\le4\cdot10^4\)，\(D\sim0.08\) 恰为应得之值。直接证实来自测试
5（高温截断：\(T_{\min}: 1000\to10000\) 时
\(D = 0.0878 \to 0.0866\)------缓慢但系统的下降）以及最重要的测试
12：真实间距与模拟 GUE 矩阵 \(2000\times2000\) 间距之间的\textbf{双样本}
KS 给出 \(D = 0.0406\)，\(p = 0.139\)------``零点间距与同尺寸 GUE
矩阵间距同分布''的假设\textbf{未被拒绝}。换言之：零点与有限尺寸 GUE
矩阵不可区分，但与抽象的无限 GUE 极限可区分------这正是有限维 GUE
相似性的定义。

\subsection{3.4 谱刚性：零点比模拟 GUE
更''刚''}\label{ux8c31ux521aux6027ux96f6ux70b9ux6bd4ux6a21ux62df-gue-ux66f4ux521a}

\begin{figure}
\centering
\includegraphics{../figures/fig03_decay_fits.png}
\caption{\textbf{图 4.} 对数坐标下 \(b(N)\) 衰减的两种拟合：幂律与
\(1/\ln N\)；残差的系统弯曲指向对数渐近。}
\end{figure}

\begin{figure}
\centering
\includegraphics{../figures/fig04_bootstrap_slope.png}
\caption{\textbf{图 5.} 斜率的自助分布（1000 次重采样）：95\% CI
\([-0.190;-0.155]\)；值 \(-0.5\) 远在区间之外。}
\end{figure}

\begin{figure}
\centering
\includegraphics{../figures/fig06_sigma2_L.png}
\caption{\textbf{图 6.} 数目方差 \(\Sigma^2(L)\)：数据（平坦）对对数型
GUE 与线性泊松------9/9 点更接近 GUE。}
\end{figure}

\begin{figure}
\centering
\includegraphics{../figures/fig07_delta3_L.png}
\caption{\textbf{图 7.} 谱刚性 \(\Delta_3(L)\)：大 \(L\) 处数据相对模拟
GUE 的超刚性。}
\end{figure}

本块信息量最大的结果来自 \(\Sigma^2(L)\) 与 \(\Delta_3(L)\)（测试
13、14）。在九点网格 \(L\in[2;1400]\) 上，零点数据给出
\(\Sigma^2 \approx 0.33\text{–}0.41\)------近乎常数；而泊松预测线性增长至
1400，模拟 GUE 参考对数增长至 \(0.61\text{–}0.72\)。数据在全部 9/9
个点上都比泊松更接近 GUE；但在大 \(L\) 处它们甚至比模拟 GUE
\textbf{更刚}（更小）------数据 \(\Delta_3\) 相对 GUE 参考的超额从
\(L=2\) 的 \(0.99\) 降到 \(L=1400\) 的 \(0.48\)。这种超刚性是真实 ζ
零点的已知特征（算术系统的谱刚性高于矩阵系统）；展析伪影只会抬高
\(\Delta_3\) 与 \(\Sigma^2\) 而非压低，因此其复现是展析质量的独立证明。

\section{4. 第二块：AB 云构造与拓扑定理（测试
15、17、20、24--27）}\label{ux7b2cux4e8cux5757ab-ux4e91ux6784ux9020ux4e0eux62d3ux6251ux5b9aux7406ux6d4bux8bd5-1517202427}

\subsection{4.1
构造的正确性：厄米性与磁通}\label{ux6784ux9020ux7684ux6b63ux786eux6027ux5384ux7c73ux6027ux4e0eux78c1ux901a}

在谈统计之前，构造必须精确。测试 15 在 \(16\times16\) 晶格、双涡旋
\(q=\pm1\) 下检查：厄米性
\(\|H-H^\dagger\|/\|H\| = 0\)（机器零）；对称类------\(\tau(H)=\|H-H^*\|/\|H\|=0\)，即狄拉克规范下矩阵为实（Dyson
三分法预言的 GOE 类）；每个涡旋元胞及对照空元胞的磁通准确到 \(10^{-14}\)
弧度。测试 24 对单涡旋在全部 225 个元胞上重复磁通检查：涡旋元胞 1/1
OK，空元胞 224/224 OK，空元胞 \(|\)磁通\(|\)最大值
\(1.47\cdot10^{-14}\)。这两项检查是地基：相位构造的任何错误都会破坏磁通，而磁通是完美的。

\subsection{4.2
Byers--Yang：整数涡旋不可见，分数涡旋是物理的}\label{byersyangux6574ux6570ux6da1ux65cbux4e0dux53efux89c1ux5206ux6570ux6da1ux65cbux662fux7269ux7406ux7684}

\begin{figure}
\centering
\includegraphics{../figures/fig14_byers_yang.png}
\caption{\textbf{图 8.} Byers--Yang
检验：整数荷在谱上不可见（\(3.5\cdot10^{-15}\)），分数荷是物理的（\(0.34\)）；对数标度。}
\end{figure}

\textbf{Byers--Yang
定理}（1961）的格点表述：谱对每元胞添加整数磁通量子周期------整数荷涡旋可被只改变相位的幺正格点变换''规范掉''。测试
25 直接测量：\(16\times16\) 开边界、\(\alpha=0\)、两个四中性涡旋
\(q=\pm1\) 随机构型相对干净晶格给出最大谱移
\(3.49\cdot10^{-15}\)------整数涡旋在谱上不可见，恒等式精确。单个分数构型
\(q=0.3\) 给出谱移 \(0.342\)------分数涡旋携带物理磁通
\(4\pi\cdot0{,}3 = 1.2\pi \not\equiv 0 \pmod{2\pi}\)，真实地扰动谱。支配后续全部物理的推论：狄拉克规范下只有\textbf{分数}荷或背景
\(\alpha\neq0\) 能产生 GUE 统计；光滑规范 \texttt{monumental}
中相位对任意 \(q\) 均为复数，整数荷亦可（第 5
章）。两个分支均已数值验证；这一分叉解释了套件各版本间的大部分历史分歧。

\subsection{4.3 Connes 自对偶与 AIII
手征对称}\label{connes-ux81eaux5bf9ux5076ux4e0e-aiii-ux624bux5f81ux5bf9ux79f0}

\(\alpha = 1/2\) 处 Hofstadter 哈密顿量具有 \textbf{Connes
自对偶}：\(\alpha\to1/\alpha\) 等价于对格点的傅里叶变换，\(\alpha=1/2\)
为不动点。环面上测试 17 测得：\(\alpha=1/2\) 时零模数恰为 4（偶 \(L\)
时为不动点棱数的两倍------解析值），E→−E 谱对称缺陷
\(1.4\cdot10^{-15}\)。\textbf{子格手征对称}（测试
27）：在二部正方晶格上，算符 \(\Gamma\)（棋盘符号）在无涡旋时对任意
\(\alpha\) 满足 \(\Gamma H\Gamma = -H\)；\(W=0\) 时测得缺陷为
\(0\)（精确），\(W=4\) 时
\(5.9\cdot10^{-3}\)（小而非零，格点无序破坏二部性），双涡旋无无序时又恰为
\(0\)：相位涡旋保持二部性。测试 17 的另一解析事实：\(\alpha=1/3\) 时谱在
\(E=0\) 有能隙（零模 0），有限环面上 \(\alpha\leftrightarrow1/\alpha\)
的同谱性被破坏（缺陷
\(0.30\)）------它仅在布洛赫带意义下精确；我们如实记录，有限尺寸的破坏并不与无穷体积定理矛盾。

\subsection{4.4 陈数：C₁ = 2}\label{ux9648ux6570cux2081-2}

测试 20 用 Fukui--Hatsugai--Suzuki 方法在 \(6\times6\)
环面、\(\alpha=1/2\)（无涡旋）下计算第一能带\textbf{陈数}：\(8\times8\)
扭转边界条件网格，最低带的 Berry 截面给出 \(C_1 = 2\)。专著曾预言每子带
\(|C_1|=1\)；\(\alpha=1/2\) 处 Hofstadter 带成对合并，``最低带''含两个
\(C_1=\pm1\) 贡献同号的子带------正确计入简并后总 \(|C_1|=2\)
与预言一致。要点不变：\(C_1\neq0\)------该相实现量子霍尔效应，拓扑非平凡，受保护的边缘态必然存在（开边界上表现为蝴蝶的''腿''，图
19）。

\subsection{4.5
环面与时间反演破缺}\label{ux73afux9762ux4e0eux65f6ux95f4ux53cdux6f14ux7834ux7f3a}

测试 26 在周期晶格上合并磁通与统计检查：环面 \(16\times16\)，单涡旋
\(q=1\) 满足规范相容条件 \(q(2-N_x) = -14\in\mathbb{Z}\)，沿 +x
的周期狄拉克弦。磁通：涡旋元胞与全部 255
个空元胞------OK；厄米性------OK；用于统计的光滑规范下
\(\tau(H) = 0.2252\)------T 对称破缺（GUE 类），且
\(\langle r\rangle = 0.5857\)，位于 GUE \(0.5992\)
校准带内。该测试是整个架构的缩影：单一构型中同时具备精确磁通 + T 破缺 +
GUE 统计。

\section{5. 第三块：AB 云的 GUE 分类（测试
16、26、29、33、34）}\label{ux7b2cux4e09ux5757ab-ux4e91ux7684-gue-ux5206ux7c7bux6d4bux8bd5-1626293334}

\subsection{5.1 规范的角色：同一模型为何给出 GOE 与
GUE}\label{ux89c4ux8303ux7684ux89d2ux8272ux540cux4e00ux6a21ux578bux4e3aux4f55ux7ed9ux51fa-goe-ux4e0e-gue}

测试 16（Montgomery
检验）在狄拉克规范、\(N_v=2\)、\(q=\pm1\)、\(\alpha=1/2\) 下给出
\(\langle r\rangle = 0.5133\pm0.0353\)------介于 GOE（0.531）与
GUE（0.600）之间，对 30 个 \(256\times256\) GUE 矩阵的蒙特卡洛 p 值为
\(0.00\)。原因见
4.2：该规范下整数荷规范不可见，矩阵近乎实数，GOE''天花板''无法被任何无序克服------这是矩阵层面的、而非噪声层面的
T 对称。测试 16 在测试 33 内以光滑规范重复：相同涡旋位置下
\(\langle r\rangle\) 升至
\(0.585\)。对可复现性的教训：\textbf{相位规范是物理规约的一部分}；不指明规范地比较套件版本毫无意义------现在每份报告的标题都注明规范。

\subsection{5.2 多实现自助法（测试
33）}\label{ux591aux5b9eux73b0ux81eaux52a9ux6cd5ux6d4bux8bd5-33}

\begin{figure}
\centering
\includegraphics{../figures/fig11_r_bootstrap.png}
\caption{\textbf{图 9.} 测试 33：\(q=1\)（环面）与
\(q=0.3\)（开边界）的三次 \(\langle r\rangle\) 实现，背景为 GUE
与泊松线。}
\end{figure}

本工作的主要统计结果。构型：环面 \(16\times16\)、双涡旋
\(q=\pm1\)、\(\alpha=1/2\)、无序
\(W=4\)、光滑规范、三次独立实现（随机涡旋位置、不同无序种子）。结果：\(\langle r\rangle = 0.5639,\ 0.5811,\ 0.6094\)；合并
\(\langle r\rangle = 0.5848\pm0.0260\)，95\% CI \([0.5588;0.6108]\)，对
GUE 目标 \(0.5992\) 的偏差
\(-0.0144\)（\(-2.4\%\)）------位于校准带内，PASS。开边界分数荷
\(q=\pm0.3\)
的对照系列：\(0.4025,\ 0.4967,\ 0.4260 \Rightarrow 0.4417\pm0.0555\)------显著压低，由边缘态（泊松成分）与环面规范相容性缺失解释。两系列之差（\(0.585\)
对 \(0.442\)）是边界条件影响有限晶格统计的最直观演示（图 11）。

\subsection{5.3 随晶格尺寸的标度（测试
34）}\label{ux968fux6676ux683cux5c3aux5bf8ux7684ux6807ux5ea6ux6d4bux8bd5-34}

\begin{figure}
\centering
\includegraphics{../figures/fig10_r_L_scaling.png}
\caption{\textbf{图 10.} 测试 34：\(\langle r\rangle\)
随晶格尺寸的收敛；两支都落到 GUE 平台 \(\approx0.60\)。}
\end{figure}

\(L = 10, 20, 30, 50\) 的 \(\langle r\rangle(L)\)
标度（每点两次实现，无序 \(W=4\)）：分数支
\(q=0.3\)（开边界）------\(0.5482 \to 0.5921 \to 0.6048 \to 0.5953\)；整数支
\(q=1\)（环面）------\(0.5739 \to 0.5681 \to 0.6042 \to 0.5864\)。两支均在
\(L\ge30\) 进入平台 \(0.59\)--\(0.605\)，在 GUE \(0.5996\)
的两个标准误之内；分数支单调，其平台与专著预言（\(0.487\to0.562\to0.595\to0.594\)）在收敛形状上吻合------我们的
\(L=10\) 值更高（我们的涡旋构型在小 \(L\)
时边缘效应更弱），此后趋势一致。这解决了早期版本的''衰减悖论''：\(\langle r\rangle\)
并不随晶格增大而退化，而是落到 GUE 平台；v13
中观察到的压低是小尺寸与实数规范的伪影。

\subsection{5.4 汇总量 f\_GUE（测试
29）}\label{ux6c47ux603bux91cf-f_gueux6d4bux8bd5-29}

测试 29
把统计压缩到单一标度：\(16\times16\)、\(N_v=2\)、\(q=\pm1\)、\(\alpha=1/2\)、\(W=4\)
时，\(\Sigma^2(L{=}2) = 0.8377\) 给出 \(f_{\rm GUE} = 0.648\)，PASS 阈值
\(0.90\)------WARN。这与测试 33 并不矛盾：小 \(L\) 的 \(\Sigma^2\)
对谱的刚性核心与边缘敏感，那里混合尚未普适；\(0.90\) 阈值是按
\(\ge100\times100\) 系综校准的。第 4 块中我们将在 \(K(t)\)
上看到同一图景：小系综显示正确的\emph{趋势}但达不到按更大系统校准的硬阈值。我们认为这种指标分歧不是失败，而是有限尺寸的可测特征，已记入汇总表。

\section{6. 第四块：与 ζ 零点的直接比较（测试
28、35、36）}\label{ux7b2cux56dbux5757ux4e0e-ux3b6-ux96f6ux70b9ux7684ux76f4ux63a5ux6bd4ux8f83ux6d4bux8bd5-283536}

\subsection{6.1 零点样本的双通比较（测试
28）}\label{ux96f6ux70b9ux6837ux672cux7684ux53ccux901aux6bd4ux8f83ux6d4bux8bd5-28}

\begin{figure}
\centering
\includegraphics{../figures/fig15_berry_R2_cutoff.png}
\caption{\textbf{图 11.} Berry 修正 \(R_2(0;T)\)：有限样本截断与向
\(-1\) 渐近的缓慢逼近；竖线标记两通高度（1000/5000 零点）。}
\end{figure}

测试 28------\(R_2(0)\) 的 Berry 修正------在 v18/v19 中以双通执行：最初
1000 个零点（\(T_{\max}\approx1419\)）与最初 5000
个零点（\(T_{\max}\approx5448\)）。素数和
\(\sum_p\sum_{k\le5}p^{-k/2}/(k\ln p) = 4.2079\) 给出修正后的理论值
\(R_2(0) = -2.3394\)；预期的有限样本附加量 \(\sim T^{-1/2}\) 从
\(0.0019\)（\(T=1419\)）降到
\(0.0009\)（\(T=5448\)）。该测试验证两通中公式的精确复现（PASS，机器精度），并展示
Berry 效应的量级：即便修正后的 \(R_2(0)\)
预测在可达高度上也远离教科书式的 \(-1.007\)
渐近值。对套件的实践意义：两通设定了解释任何小零点样本''吻合''所需的标度。

\subsection{6.2 Montgomery 检验：AB 云对 ζ-5000（测试
35）}\label{montgomery-ux68c0ux9a8cab-ux4e91ux5bf9-ux3b6-5000ux6d4bux8bd5-35}

\begin{figure}
\centering
\includegraphics{../figures/fig08_r2_pair.png}
\caption{\textbf{图 12.} 对关联 \(R_2(s)\)：AB 云对 ζ 零点、GUE
与泊松------两支都复现关联洞。}
\end{figure}

正面对比：AB 云展析间距（2 次实现，304 个间隙）对最初 5000 个 ζ
零点的展析间距（4999 个间隙，对称滑窗展析）。结果：双样本 KS
\(D = 0.1173\)，\(p = 6.7\cdot10^{-4}\)------分布完全相同的假设被拒；对关联差的平均模
\(\langle|R_2^{\rm AB}(s)-R_2^{\zeta}(s)|\rangle = 0.147\)，阈值
\(0.10\)。但同一张 \(R_2(s)\) 表（\(s\) 从 0.05 到 2.85 的 29 个值；图
8）显示要点：两者都有关联洞------\(R_2^{\rm AB}(0.05) = 0.00\)，\(R_2^{\zeta}(0.05) = 0.005\)，泊松为
1；两条曲线平行于 GUE 预测，且都通过''比泊松更接近
GUE''的检查（\(d_{\rm GUE} = 0.140 < d_{\rm Pois} = 0.227\)）。诚实的总结：\textbf{AB
云复现了零点关联的定性与半定量结构（洞、\(s\approx1\)
处的折点、平台），但并不逐点吻合}------304 个间隙不足以支撑稳定的
\(R_2\) 或 KS
级别的吻合。我们把''同一普适类''与''同一谱''之间的距离记为数字，而非形容词。

\subsection{6.3 形因子 K(t)（测试
36）}\label{ux5f62ux56e0ux5b50-ktux6d4bux8bd5-36}

\begin{figure}
\centering
\includegraphics{../figures/fig09_K_form_factor.png}
\caption{\textbf{图 13.} 谱形因子 \(K(t)\)：GUE 斜坡+平台、AB 云数据与 ζ
零点；小系综显示正确趋势但带离群点。}
\end{figure}

谱形因子是对要求最高的诊断：稳定估计需要数百个独立间隙与精确展析。我们的
306 个本征值给出对 GUE 斜坡的 RMS 偏差 \(0.93\)（PASS 阈值
\(0.30\)）、与理论曲线的相关 \(0.018\)（阈值 \(0.50\)）------两项均为
WARN；曲线形状（图 9）显示向 \(t\approx1\)
的正确上升及随后的回落，带有小系综典型的离群点。同一图上的 ζ
零点曲线（取自 5000 个零点）对 GUE 的 RMS 为
\(0.86\)------即使参考零点在该尺度也过不了硬阈值：该诊断有信息量，但按大若干量级的系综校准。我们把它纳入专著，作为适用性界限的演示与未来大运行的课题。

\section{7. 第五块：狄拉克动力学与非厄米性（测试
19、30--32、37）}\label{ux7b2cux4e94ux5757ux72c4ux62c9ux514bux52a8ux529bux5b66ux4e0eux975eux5384ux7c73ux6027ux6d4bux8bd5-19303237}

\subsection{7.1 狄拉克锥（测试
19、31）}\label{ux72c4ux62c9ux514bux9525ux6d4bux8bd5-1931}

\begin{figure}
\centering
\includegraphics{../figures/fig12_dirac_cone.png}
\caption{\textbf{图 14.}
狄拉克锥：\(E_{\min} = 0.02 + 11.83/L\)，\(R^2 = 0.9997\)------每点均有四个零模时对
\(1/L\) 的线性。}
\end{figure}

\(\alpha=1/2\) 时 AB
云中的涡旋表现得像一维相对论费米子：零附近的谱是线性的。数值检查（测试
19）在 \(L = 12, 16, 24, 32, 48\) 晶格上运行：最小能量
\(E_{\min}\)（四个零模下 E=0 处的能隙）按
\(E_{\min} = 0.0202 + 11.826/L\) 标度，\(R^2 = 0.9997\)------对 \(1/L\)
的线性无可挑剔。常数的解读：若狄拉克点位于布里渊区中心（\(k_{\min} = 2\pi/L\)），有效费米速度
\(v_F = b/2\pi \approx 1.88\)（以晶格常数 \(b\)
上的跃迁为单位）；若在区边缘（\(k_{\min}=\pi/L\)），\(v_F = b/\pi \approx 3.76\)。历史值
\(\sim0.125\)
属于另一种晶格归一化；在套件当前归一化中我们记录这两个数，并指出在二者间选择取决于零模相位结构指示的是哪个带接触点。测试
31（周期边界的环面变体）显示 \(R^2 = 0.03\)
且能隙为零：环面上零模受拓扑保护，任何 \(L\)
都不会打开------这本身就是拓扑本性的确证，但能隙标度需要开边界。

\subsection{7.2 态密度的狄拉克凹陷（测试
30）}\label{ux6001ux5bc6ux5ea6ux7684ux72c4ux62c9ux514bux51f9ux9677ux6d4bux8bd5-30}

\begin{figure}
\centering
\includegraphics{../figures/fig13_dirac_dip.png}
\caption{\textbf{图 15.} \(\alpha=1/2\) 处 DOS 的狄拉克凹陷：双向对比度
\(20\times\)（对数标度）。}
\end{figure}

\(\alpha=1/2\) 处 Hofstadter 谱有解析特征：\(E=0\)
附近态密度线性归零（狄拉克点），而 \(\alpha\neq1/2\) 时谱是普通的。测试
30 在 \(20\times20\) 无无序晶格上测量中央带 \(|E|<0.5\)
的态密度：\(\rho(\alpha=0.4) = 0.200\)，\(\rho(\alpha=0.5) = 0.010\)，\(\rho(\alpha=0.6) = 0.200\)------双向对比度恰为
\(20\times\)。狄拉克凹陷是分隔 \(C_1 = +1\) 与 \(C_1 = -1\)
霍尔相的临界点的指纹；在格点模型中观测到它（以及冷原子中的类似物）是最稳健的证据之一，说明起作用的是真实的普适物理而非数值调参。

\subsection{\texorpdfstring{7.3 Hatano--Nelson：离开
\(\sigma=1/2\)（测试
32）}{7.3 Hatano--Nelson：离开 \textbackslash sigma=1/2（测试 32）}}\label{hatanonelsonux79bbux5f00-sigma12ux6d4bux8bd5-32}

\begin{figure}
\centering
\includegraphics{../figures/fig16_hatano_nelson.png}
\caption{\textbf{图 16.}
Hatano--Nelson：\(\sigma=0.5\)------实谱；\(\sigma=0.7\)------整个谱进入复平面（趋肤效应）。}
\end{figure}

广义模型中的参数 \(\sigma\) 设定跃迁相位的''分数''部分；\(\sigma=1/2\)
时模型厄米。测试 32 在 \(16\times16\) 上比较 \(\sigma=0.5\) 与
\(\sigma=0.7\)：\(\sigma=0.5\) 时最大
\(|\mathrm{Im}\,E| = 0\)（谱严格为实）；\(\sigma=0.7\) 时全部 256
个能级为复数，最大 \(|\mathrm{Im}\,E| = 0.623\)，Hatano--Nelson 趋肤参数
\(g = (\sigma-1/2)\ln|\gamma| = 0.322\)。含义：\textbf{离开
\(\sigma=1/2\)
把系统从幺正（GUE）类移入耗散（无关联）类；本征矢向边界的局域化（趋肤效应）摧毁体统计，非厄米支上
\(\langle r\rangle\) 塌缩为零（所有间隙在复平面上''配对''）}。临界线
\(\sigma=1/2\) 不仅由 Connes
自对偶凸显，而且是存在体随机矩阵统计的唯一线条（图
13、16）。这是''临界线最优性''的第二个独立论证------解析自对偶的统计对应。

\subsection{7.4 字节级稳健性（测试
37）}\label{ux5b57ux8282ux7ea7ux7a33ux5065ux6027ux6d4bux8bd5-37}

可复现性的实践问题：结论能否经受量子化？测试 37 验证
\(\langle r\rangle\) 在本征值 256 级量子化下存活（变化 \(0.0015\)，容差
\(0.02\)），且生成器字节流的香农熵为 \(7.9991\) 比特（最大 8；基于 CLT
的 \(\chi^2\)
判据：\(z=-0.234\)，\(|z|<3\)）。该测试还记录了全局与局部分类之差：特定实现的全局谱看似泊松（\(\langle r\rangle_{\rm raw}=0.447\)），但中央
15 能级窗口内 \(\langle r\rangle_{\rm local}=0.549\)------即 5.4
节已述的体--核之别。

\section{8. 第六块：解析诠释（测试 21--23
及延伸）}\label{ux7b2cux516dux5757ux89e3ux6790ux8be0ux91caux6d4bux8bd5-2123-ux53caux5ef6ux4f38}

本块是从测量通向解析结构的桥梁。我们审慎区分：（a）机器精度可复现的精确公式（测试
21--23）；（b）与数字一致但不能由套件严格推出的物理诠释；（c）作为类比呈现的类比。这种三层呈现是方法论原则。

\subsection{\texorpdfstring{8.1 旋量相位 \(\gamma^{\ast}\)：经由虚单位的
90°（测试
21）}{8.1 旋量相位 \textbackslash gamma\^{}\{\textbackslash ast\}：经由虚单位的 90°（测试 21）}}\label{ux65cbux91cfux76f8ux4f4d-gammaastux7ecfux7531ux865aux5355ux4f4dux7684-90ux6d4bux8bd5-21}

\begin{figure}
\centering
\includegraphics{../figures/fig18_gamma_phase.png}
\caption{\textbf{图 17.} 旋量相位 \(\gamma^{\ast} = a_C + i b_C\)：辐角
\(89.87°\) 对参考 \(90°\)；\(b_C\) 的主导把相位压向 \(\pi/2\)。}
\end{figure}

定义 \(\delta_C = \pi/7 = 0.448799\) 弧度（见 8.2）与
\(b_2(K3) = 22\)（K3 曲面的第二 Betti 数）。测试 21
验证两个实数量：\textbf{实制动项}
\(a_C(\gamma^{\ast}) = \delta_C^5/b_2(K3) = 8.276\cdot10^{-4}\)
与\textbf{虚 Berry 项}
\(b_C(\gamma^{\ast}) = 1-\cos(2\delta_C) = 2\sin^2\delta_C = 0.376510\)；两者均以机器精度复现。复旋量相位
\(\gamma^{\ast} = a_C + i b_C\) 的模为
\(|\gamma^{\ast}| = 0.3765\)，辐角 \(89.874°\)------对 \(\pi/2\)
的偏离仅 \(0.126°\)（\(2.2\cdot10^{-3}\) 弧度）。结构如下：由于
\(b_C \gg a_C\)，\(\gamma^{\ast}\) 的辐角自动压向 90°------``90°
的跃变经由虚单位实现''：作为几何对象的相位被构造得使其相对论（旋量）本性以纯虚优势表达。把它诠释为假想的
Peccei--Quinn 相位属于（b）层：与数字相容，但不能由套件推出。

\subsection{\texorpdfstring{8.2 AB 相位 \(\Phi_{AB} = \pi/7\)
与分数电荷（测试
22）}{8.2 AB 相位 \textbackslash Phi\_\{AB\} = \textbackslash pi/7 与分数电荷（测试 22）}}\label{ab-ux76f8ux4f4d-phi_ab-pi7-ux4e0eux5206ux6570ux7535ux8377ux6d4bux8bd5-22}

精确公式：荷电 \(q/e = 1/14\)
的粒子环绕一个磁通量子的通量管时，阿哈罗诺夫--玻姆相位为
\(\Phi_{AB} = 2\pi/14 = \pi/7 = 0.448799\) 弧度 \(= 25.714°\)。测试 22
在精确有理数算术中确认该恒等式（偏差 \(0.00\)）。量 \(\delta_C=\pi/7\)
正是 \(\Phi_{AB}\)；也就是说 \(q/e=1/14\)
是分数电荷假说，整个解析上层建筑围绕它建立。值得注意的是，\(\pi/7\) 与
PSL(2,7) 特征表的的高斯和 \(b_7=(-1+i\sqrt7)/2\) 相联系（8.5
节）：七次单位根既出现在 Klein 四次曲面的旋量结构中，也出现在 AB
相位中------这是 v21 专著已注意到的结构押韵。

\subsection{8.3 复分形因子 CF 与对数抛物化（测试
23）}\label{ux590dux5206ux5f62ux56e0ux5b50-cf-ux4e0eux5bf9ux6570ux629bux7269ux5316ux6d4bux8bd5-23}

对具有离散标度比 \(\lambda_{\rm DSI} = b_2(K3) = 22\) 的临界解，测试 23
计算\textbf{复分形因子}

\[CF = \exp\!\left(\frac{2\pi i}{\lambda_{\rm DSI}}\right)\cdot\exp\!\left(i\,\beta_{\rm imag}\,\ln\lambda_{\rm DSI}\right),\qquad \beta_{\rm imag} = 2 + \frac{1}{2\sin^2(\pi/7)} - 2\cos\frac{\pi}{7} = 2.854033,\]

给出 \(CF = -0.9501 + 0.3119\,i\)，按构造
\(|CF| = 1\)。注意指数内部的对数------这就是''对数抛物化''：用对数相位调制离散标度不变性。测试还带有诚实的告示：晚期
tex 源中的硬编码值 \(CF_{\rm hard} = 0.9786 + 0.0390\,i\)（模
\(0.9794\)，由此
\(c_{AB} = 1-|CF_{\rm hard}| = 0.0206\)）是\textbf{经验输入}------类似
QCD 中的 \(\alpha_s\)------而公式是结构（拓扑）预测；两幅度恰差
\(c_{AB}\)，测试连同各自状态记录两个数。

\subsection{8.4 T 破缺作为 GUE
之源：黑洞类比}\label{t-ux7834ux7f3aux4f5cux4e3a-gue-ux4e4bux6e90ux9ed1ux6d1eux7c7bux6bd4}

第 3--5 块的数字拼合成 Dyson 的解析图景：GOE ↔ T 不变性，GUE ↔
其破缺。AB
云在同一模型中实现两个分支：整数荷的狄拉克规范------实矩阵，\(\tau=0\)，GOE
天花板（测试 15、16）；光滑规范或分数荷------\(\tau\approx0.225\)，GUE
统计（测试 26、33）。讨论中提出的假设强化了这条线：\textbf{T
破缺与黑洞统计由幺正系综描述是同一原因（那里时间被弯曲，热力学上反转不是对称）}。我们强调该构造的地位：它不是套件的推论，而是套件不与之矛盾的诠释框架------并且它使预言''任何具有复跃迁且无可积性的相位模型都将给出
\(\langle r\rangle\to0.6\)``可以在独立材料上检验。

\subsection{8.5 Kerr 类比、能层与 PSL(2,7) ⊂
Monster（综述）}\label{kerr-ux7c7bux6bd4ux80fdux5c42ux4e0e-psl27-monsterux7efcux8ff0}

v21 的两个诠释以综述地位保留。\textbf{Kerr
类比}：涡旋相位场定义''能层区''------相位速度相对晶格变号的区域，如 Kerr
能层；其数值验证需要动力学而非静态谱，仍是课题。\textbf{Monster
桥}：维数 196883 的 Monster 特征 \(\chi_2\) 限制到子群
PSL(2,7)（经极大子群 \(2^{3+6+12+18}.(L_3(2)\times3S_6)\)
嵌入）近似分解为
\(1456\cdot\chi_1 + 3334\cdot\chi_3 + 3334\cdot\chi_3' + 6784\cdot\chi_6 + 8081\cdot\chi_7 + 9770\cdot\chi_8\)（维数和
196891，偏差 0.004\%）；重数可解释为多涡旋态的选择定则。精确分解需要 GAP
CTblLib 的 class
fusion------开放课题；我们在任何裁决中都不使用近似系数。

\section{9. 综合：作为容许态编码的 ζ
零点}\label{ux7efcux5408ux4f5cux4e3aux5bb9ux8bb8ux6001ux7f16ux7801ux7684-ux3b6-ux96f6ux70b9}

\begin{figure}
\centering
\includegraphics{../figures/fig19_vortex_texture.png}
\caption{\textbf{图 18.} 涡旋对 \(q=\pm1\) 叠加 Landau
规范的相位织构：产生观测统计的相位''云''。}
\end{figure}

\begin{figure}
\centering
\includegraphics{../figures/fig20_hofstadter_butterfly.png}
\caption{\textbf{图 19.} Hofstadter 蝴蝶 \(E(\alpha)\)；竖线标记自对偶点
\(\alpha=1/2\)------专著关键结果的汇聚处。}
\end{figure}

把链条拼起来。（1）光滑规范下 \(\alpha=1/2\) 的 AB 云谱具有随尺寸收敛的
GUE 统计：\(L: 10\to50\) 时 \(\langle r\rangle: 0.548\to0.595\)（测试
33、34）。（2）ζ 零点的统计是 GUE 加可计算的有限样本修正（测试
4--14）；在可达高度上''纯 GUE''被拒，``有限尺寸 GUE''不被拒（测试
12）。（3）直接比较显示 \(R_2\) 的定性结构吻合（测试
35）------Montgomery 关联洞被涡旋云复现。（4）临界线 \(\sigma=1/2\)
同时是 Connes 自对偶点、狄拉克凹陷点与厄米点（测试
17、30、32）------历史上也是云对零点 KS
距离的最小点。（5）涡旋表现得像相对论性粒子（狄拉克锥，测试
19），分数电荷充当物理磁通（Byers--Yang，测试 25）。最终叙事------``ζ
零点作为相位谐振腔容许能量态的编码''------由此获得\textbf{自洽且经定量检验的纲领}的地位，而非被证明的定理：我们清楚哪些环节机器精确（磁通、恒等式），哪些统计稳健（标度、自助），哪些仍然开放（\(T\to\infty\)
渐近、\(R_2\) 的逐点吻合）。

\section{10.
验证方法学：``PASS''究竟保证什么}\label{ux9a8cux8bc1ux65b9ux6cd5ux5b66passux7a76ux7adfux4fddux8bc1ux4ec0ux4e48}

套件 v19
是第一个在架构上把报告引擎与计算核心隔离的版本。这一决策的历史颇有教益：v18
中 PDF 写出器的一个缺陷（行列表被声明为 \texttt{String{[}{]}}
而循环推入的是元组）在\textbf{每项}测试完成后触发
\texttt{MethodError}；异常被外层处理器捕获，裁决被降级为
WARN，双通模式的第二通根本不会运行------``快速检查''退化为单通，而外表上似乎正常。一个
bug 造成的三个数字：37 项测试虚假 WARN、零份 PDF
报告、零次第二通。修复（FIX-R1\ldots R5，附录 A）确立了：（1）PDF
写出器正确；（2）任何报告格式的失败不能影响裁决；（3）测试
28（双通循环的变量作用域）与测试 18（\(\alpha=1/3\)
时环面规范相容性）已修复；（4）快速检查截断至 5000 零点并真正执行双通
16×16→32×32。教训：\textbf{套件的结论不会强于其最弱的基础设施}，而''近乎正常''的报告管线比损坏的更危险，因为它伪造了已验证的表象。

当前版本的实际保证：每项测试是带完整日志的独立目录；本书每个数字可追溯到测试编号；裁决可由形如
\texttt{julia\ ab\_cloud\_v19.jl\ -\/-test\ N\ -\/-no-two-pass}
的命令复现；完整运行与 \texttt{-\/-quick} 默认双通；凡需随机性之处 RNG
均用固定种子（逐位可复现）。

\section{11. 结论}\label{ux7ed3ux8bba}

\begin{enumerate}
\def\labelenumi{\arabic{enumi}.}
\tightlist
\item
  \textbf{构造精确性}（机器级）：厄米性 0；狄拉克弦磁通准确到
  \(10^{-14}\)；整数荷 Byers--Yang \(3.5\cdot10^{-15}\)；Connes
  自对偶------4 个零模且 E→−E 缺陷 \(10^{-15}\)；\(W=0\) 手征性
  \(0.0\)。构造在数学上正确------这是被证明的，而非估计的。
\item
  \textbf{GUE 普适性}（统计级）：\(\langle r\rangle = 0.5848\pm0.0260\)
  对 GUE \(0.5992\)（偏差 \(-2.4\%\)）；标度 \(L=10\to50\) 单调并在
  \(0.595\pm0.008\) 进入平台；谱在所有尺度 \(L\) 上比泊松更刚，9/9
  点上更接近 GUE。T 对称恰在相位为复数处破缺------Dyson 三分法得到复现。
\item
  \textbf{与 ζ 零点的比较}（诚实级）：\(R_2(s)\)
  的定性结构吻合（洞、平台；比泊松更接近 GUE），304
  个间隙下不存在也不应期待逐点吻合；Berry 有限样本修从定量上解释了对渐近
  GUE 的一切观测偏离。
\item
  \textbf{拓扑与动力学}：\(C_1=2\)（非平凡 IQHE）、狄拉克锥
  \(R^2=0.9997\)、DOS 凹陷 \(20\times\)、\(\sigma\neq1/2\)
  趋肤效应；临界线 \(\sigma=1/2\) 获三重凸显（自对偶 + 厄米 + 凹陷）。
\item
  \textbf{解析诠释}：\(\gamma^{\ast}\)、\(\Phi_{AB}=\pi/7\)、\(CF\)
  诸公式以机器精度复现并构成自洽的
  \(\pi/7\)-结构；其物理解读（Peccei--Quinn、Moonshine）仍是假说------套件不矛盾，但也不证明。
\item
  \textbf{信任的基础设施}：v19
  修复了报告管线的一切已发现缺陷；每个结论可追溯至日志；两套裁决系统（机器与校准）被明确分开。
\end{enumerate}

\section{12. 开放问题}\label{ux5f00ux653eux95eeux9898}

\begin{enumerate}
\def\labelenumi{\arabic{enumi}.}
\tightlist
\item
  \textbf{\(b(N)\)
  的渐近}：幂律还是对数？替代拟合更好（\(R^2=0.9999\)），但 \(N\)
  范围不足；需要到 \(N\sim10^6\) 的运行（加载器已支持 2M Odlyzko
  数据集）。
\item
  \textbf{\(R_2\) 的逐点吻合}：AB 云系综须增至 \(\ge10^4\)
  间隙（更大晶格 × 多实现 × 涡旋移位）------只有那时对 ζ-5000 的双样本
  KS 才有信息量。
\item
  \textbf{形因子 \(K(t)\)}：同一目标；按我们估计 PASS
  阈值（\(\mathrm{RMS}<0.30\)、\(\mathrm{corr}>0.50\)）在 \(\ge3000\)
  间隙时可达。
\item
  \textbf{3D 接口}：Hofstadter 堆叠中的涡旋线（3D 实验室模块）------GUE
  平台是否沿第三维保持；已有初步数据但无 37 测试验证。
\item
  \textbf{196883 对 PSL(2,7) 的精确分解}：需要经 GAP CTblLib 的 class
  fusion；在此之前系数仍是近似的。
\item
  \textbf{涡旋动力学与 Kerr 类比}：静态谱无法检验能层诠释；需要带 AB
  相位的点涡旋拉格朗日模拟。
\end{enumerate}

\section{附录 A.
代码：架构与关键模块}\label{ux9644ux5f55-a.-ux4ee3ux7801ux67b6ux6784ux4e0eux5173ux952eux6a21ux5757}

验证套件的完整代码（文件 \texttt{ab\_cloud\_v19.jl}，Julia
1.10，零外部依赖------包括自研的带 pHYs 的 PNG 写出器、矢量 PDF、带 LZW
的 GIF89a、stored-ZIP 的 DOCX 及全部统计）随专著发布于 \texttt{code/}
目录；无参数运行时打开交互菜单（键
\texttt{f}------快速检查）。以下列出架构模块并引用关键代码。

\textbf{文件模块}：（1）\texttt{Config}
结构与双语消息；（2）数据加载器------5 万内置 Odlyzko 零点 + 至 2M
的外部文件与优雅回退；（3）统计核心（展析、KS、χ²、AD、自助、\(\Sigma^2\)、\(\Delta_3\)、\(R_2\)、\(K(t)\)）；（4）AB
云构造器 \texttt{build\_ab\_cloud\_hamiltonian}（相位模型
\texttt{:dirac}、\texttt{:monumental}）与磁通检查
\texttt{plaquette\_flux}；（5）37 项测试 \texttt{test*}；（6）双通运行器
\texttt{run\_test\_two\_pass}；（7）v16/v19
报告引擎（\texttt{rep\_*}，md/html/pdf/docx/svg/png/gif
写出器，index.html）；（8）物理实验室（22 个实验构型）与 3D 实验室（30
项测试）；（9）菜单与 CLI。

\textbf{关键函数------哈密顿量构造器}（节选；全文见随附文件）：

\begin{Shaded}
\begin{Highlighting}[]
\CommentTok{\# φ\_ij = 2π·α·j + Σ\_k q\_k·[arg(r\_i−r\_k) − arg(r\_j−r\_k)] + φ\^{}弦\_ij}
\CommentTok{\# 键 (i→j) 取 H\_ij = −exp(iφ\_ij)；含涡旋 k 的元胞携带磁通 2π·q\_k}
\KeywordTok{function} \FunctionTok{build\_ab\_cloud\_hamiltonian}\NormalTok{(cfg}\OperatorTok{::}\DataTypeTok{ABCloudConfig}\NormalTok{)}
\NormalTok{    N }\OperatorTok{=}\NormalTok{ cfg.Nx }\OperatorTok{*}\NormalTok{ cfg.Ny}
\NormalTok{    H }\OperatorTok{=} \FunctionTok{zeros}\NormalTok{(}\DataTypeTok{ComplexF64}\NormalTok{, N, N)}
    \ControlFlowTok{for}\NormalTok{ iy }\KeywordTok{in} \FloatTok{0}\OperatorTok{:}\NormalTok{(cfg.Ny}\OperatorTok{{-}}\FloatTok{1}\NormalTok{), ix }\KeywordTok{in} \FloatTok{0}\OperatorTok{:}\NormalTok{(cfg.Nx}\OperatorTok{{-}}\FloatTok{1}\NormalTok{)}
\NormalTok{        i }\OperatorTok{=}\NormalTok{ iy }\OperatorTok{*}\NormalTok{ cfg.Nx }\OperatorTok{+}\NormalTok{ ix}
        \ControlFlowTok{for}\NormalTok{ (dx, dy) }\KeywordTok{in}\NormalTok{ ((}\FloatTok{1}\NormalTok{,}\FloatTok{0}\NormalTok{),(}\FloatTok{0}\NormalTok{,}\FloatTok{1}\NormalTok{))            }\CommentTok{\# 最近邻，两种取向}
\NormalTok{            jx, jy }\OperatorTok{=}\NormalTok{ (ix}\OperatorTok{+}\NormalTok{dx) }\OperatorTok{\%}\NormalTok{ cfg.Nx, (iy}\OperatorTok{+}\NormalTok{dy) }\OperatorTok{\%}\NormalTok{ cfg.Ny}
\NormalTok{            j }\OperatorTok{=}\NormalTok{ jy }\OperatorTok{*}\NormalTok{ cfg.Nx }\OperatorTok{+}\NormalTok{ jx}
            \ConstantTok{φ} \OperatorTok{=} \FloatTok{2}\NormalTok{π }\OperatorTok{*}\NormalTok{ cfg.alpha }\OperatorTok{*}\NormalTok{ jy              }\CommentTok{\# Landau 规范（列相位）}
            \ControlFlowTok{for}\NormalTok{ v }\KeywordTok{in}\NormalTok{ cfg.vortices                }\CommentTok{\# 涡旋相位（依模型）}
                \ConstantTok{φ} \OperatorTok{+=} \FunctionTok{vortex\_phase}\NormalTok{(v, ix, iy, jx, jy, cfg.model)}
            \ControlFlowTok{end}
\NormalTok{            H[i}\OperatorTok{+}\FloatTok{1}\NormalTok{, j}\OperatorTok{+}\FloatTok{1}\NormalTok{] }\OperatorTok{{-}=} \FunctionTok{exp}\NormalTok{(}\ConstantTok{im} \OperatorTok{*} \ConstantTok{φ}\NormalTok{)}
\NormalTok{            H[j}\OperatorTok{+}\FloatTok{1}\NormalTok{, i}\OperatorTok{+}\FloatTok{1}\NormalTok{] }\OperatorTok{{-}=} \FunctionTok{exp}\NormalTok{(}\OperatorTok{{-}}\ConstantTok{im} \OperatorTok{*} \ConstantTok{φ}\NormalTok{)          }\CommentTok{\# 厄米共轭}
        \ControlFlowTok{end}
\NormalTok{        H[i}\OperatorTok{+}\FloatTok{1}\NormalTok{, i}\OperatorTok{+}\FloatTok{1}\NormalTok{] }\OperatorTok{+=} \FunctionTok{disorder\_epsilon}\NormalTok{(ix, iy, cfg)   }\CommentTok{\# ε \textasciitilde{} U(−W/2, W/2)，固定种子}
    \ControlFlowTok{end}
    \ControlFlowTok{return}\NormalTok{ H}
\KeywordTok{end}
\end{Highlighting}
\end{Shaded}

\textbf{元胞磁通检查}（狄拉克规范）：含涡旋元胞四键相位之和给出
\(2\pi q \pmod{2\pi}\)，空元胞为 0；由
\texttt{plaquette\_flux(H,\ Nx,\ Ny,\ ix,\ iy)} 计算，在测试 15/24/26
中验证。

\textbf{裁决}：\texttt{log\_result}
写出裁决行与完整日志；\texttt{run\_test\_two\_pass}
在主晶格与加大晶格上运行测试，在 \texttt{finally}
中恢复配置（通隔离）；\texttt{rep\_finish!}
在写报告\textbf{之前}记录裁决，写出器故障（FIX-R2）不能改变它。

\textbf{v19 修复}（完整注释见代码）：FIX-R1------PDF
写出器中间行列表的类型标注（消除使裁决与第二通归零的大规模
MethodError）；FIX-R2------\texttt{\_write\_test\_report!} 外包
try/catch；FIX-R3/3b------测试 28 的
\texttt{last\_R2}/\texttt{berry\_ref} 置于函数作用域；FIX-R4------测试
18 中 \(\alpha\in\{1/2,1/3\}\) 时晶格跳到 6
的倍数；FIX-R5/5b------快速检查：ζ 截断至 5000、关闭 GIF、新增
\texttt{-\/-quick} CLI 标志。

\section{附录 B.
复现指南}\label{ux9644ux5f55-b.-ux590dux73b0ux6307ux5357}

\begin{Shaded}
\begin{Highlighting}[]
\CommentTok{\# 全套件（双通，需数分钟至数十分钟）：}
\ExtensionTok{julia}\NormalTok{ ab\_cloud\_v19.jl }\AttributeTok{{-}{-}test}\NormalTok{ all}
\CommentTok{\# 单项测试：}
\ExtensionTok{julia}\NormalTok{ ab\_cloud\_v19.jl }\AttributeTok{{-}{-}test}\NormalTok{ 33 }\AttributeTok{{-}{-}no{-}two{-}pass}
\CommentTok{\# 快速检查 16×16 → 32×32（ζ ≤ 5000，数分钟）：}
\ExtensionTok{julia}\NormalTok{ ab\_cloud\_v19.jl }\AttributeTok{{-}{-}quick}
\CommentTok{\# 语言：}
\ExtensionTok{julia}\NormalTok{ ab\_cloud\_v19.jl }\AttributeTok{{-}{-}lang}\NormalTok{ en }\AttributeTok{{-}{-}quick}
\CommentTok{\# 报告位置：results/run\_\textless{}日期\textgreater{}/test\_NN\_\textless{}slug\textgreater{}/\{report.md,html,pdf,docx, plots/, logs/\}}
\end{Highlighting}
\end{Shaded}

本书数字对应 2026-09-02 的完整双通运行 \textbf{v19
\texttt{run\_20260902\_134759}}（37 项测试，双通 72×72 → 96×96，含第二轮
HARDCORE 审计；测试 19/29/31 附第三轮系列；Julia 1.12.0；50,000 个
Odlyzko 零点）。第一轮汇总：\textbf{32 PASS / 5 WARN}（测试
4、5、6、11、12、28、35------统计类 WARN 判定，源于 Wigner
近似的``地板效应''，并以对精确 GUE 参照的双样本 KS
行进行校准；套件语义未变）；\textbf{第二轮 HARDCORE 全部 PASS}（共 77
条判定行）。全套工件------各测试报告（report.md/pdf/docx/html）、计算日志（logs/）、汇总
FINAL\_REPORT 与
index.html------已随仓库提供：\texttt{results/run\_20260902\_134759/}。此外，克莱因四次曲面的全部
64 个旋量结构已由参考实现 \texttt{verification/spinor64} 验证（附录
D）。早期单通 v18
基线（\texttt{results/verification\_run\_v18\_37tests\_2026-08-28.txt}）作为历史记录保留。

\section{附录 C.
图形索引}\label{ux9644ux5f55-c.-ux56feux5f62ux7d22ux5f15}

全部图形为 600 dpi PNG，标题采用相应出版语言；文件位于
\texttt{monographs/\textless{}lang\textgreater{}/figures/}。图
1------\(b(N)\) 收敛（双拟合）；图 2------ζ 间距直方图对 Wigner
猜想曲线；图 3------对数坐标下的双衰减拟合；图 4------斜率的自助分布与
95\% CI；图 5------KS \(D(T_{\min})\) 与临界线；图
6------\(\Sigma^2(L)\)；图 7------\(\Delta_3(L)\)；图
8------\(R_2(s)\)：AB 云对 ζ、GUE、泊松；图 9------形因子 \(K(t)\)；图
10------\(\langle r\rangle(L)\) 标度（\(q=0.3\) 与 \(q=1\)）；图
11------\(\langle r\rangle\) 自助法（测试 33）；图 12------狄拉克锥
\(E_{\min}(1/L)\)；图 13------DOS 狄拉克凹陷；图
14------Byers--Yang（整数对分数电荷）；图 15------Berry 截断
\(R_2(0;T)\)；图 16------Hatano--Nelson 非厄米谱椭圆；图
17------旋量相位 \(\gamma^{\ast}\)；图 18------涡旋对的相位织构；图
19------Hofstadter 蝴蝶（标注 \(\alpha=1/2\)）。

\section{附录 D（v22.1，spinor64）：克莱因四次曲面全部 64
个旋量结构的验证}\label{ux9644ux5f55-dv22.1spinor64ux514bux83b1ux56e0ux56dbux6b21ux66f2ux9762ux5168ux90e8-64-ux4e2aux65cbux91cfux7ed3ux6784ux7684ux9a8cux8bc1}

参考实现 \texttt{verification/spinor64}（Python/NumPy；复现命令：
\texttt{python3\ verification/spinor64/run\_spinor64.py}）对克莱因四次曲面的全部
64 个旋量结构 \(\varepsilon\in\mathbb{F}_2^6\) 执行两项独立实验，并修正
v21 专著中关于 idx=38 ``唯一性''的论断。

\textbf{E1 ------ 精确对称性（克莱因图 \{3,7\}）。} 以正则映射 \{3,7\}
离散化克莱因 四次曲面（56 顶点、84 边、24 个七边形，自同构群
PSL(2,7)，阶 168）。64
个自旋结构实现为具有奇面奇性的边符号（Kasteleyn/Cimasoni-Reshetikhin
模型，规范定准：生成树边均为正）。数值计算得 PSL(2,7) 作用的轨道分解
\textbf{28 / 21 / 7 / 7 / 1}；28
元轨道------奇结构（Arf=1）------为单一轨道， 证实 PSL(2,7) 在 28
条双切线上传递（Riemann--Klein 定理）。轨道内的符号
狄拉克算子互为置换共轭：全部 64 个结构的最大成对谱距为
\textbf{8.9·10⁻¹⁵}； 规范不变性 7.1·10⁻¹⁵；离散算子零模数：2（奇轨道）/
3（偶轨道）/ 7
（平凡类）。结论：任何旋量结构都不可能拥有独特统计------v21 中 ``idx=38
的唯一性''是破坏 PSL(2,7) 对称性的离散化伪影。按 v21 专著自身的公式
\(\mathrm{Arf}(\varepsilon)=\varepsilon_1\varepsilon_2+\varepsilon_3\varepsilon_4+\varepsilon_5\varepsilon_6\)，
向量 \(\varepsilon(38)=(0,1,1,0,0,1)\) 给出 Arf = 0（而非此前声称的
1）。

\textbf{E2 ------ 统计（AB 云）。} Hofstadter 环面
\(L=44\)、\(\alpha=1/2\)、 \(N_v=54\) 个 \(q=+1\) 涡旋（按密度定标锚点
\(N_v=25\)@\(30\times30\)）、\(W=0\)、 \texttt{:monumental} 规范（atan
平滑、竖向键、因子 0.5），自旋结构以边界扭转
\(\varphi_x=\pi(\varepsilon_1+\varepsilon_3+\varepsilon_5)\)、
\(\varphi_y=\pi(\varepsilon_2+\varepsilon_4+\varepsilon_6)\) 引入；对 5
组涡旋构型取平均，体窗 0.6。参照：100 个 \(1936\times1936\) GUE 矩阵的
蒙特卡罗系综（测试 16 方法论）：中位数
\(\langle r\rangle = 0.6013\)，95\% 置信区间
\([0.5847; 0.6140]\)。结果：\textbf{全部 64 个结构均落于置信区间内}：
\(\langle r\rangle = 0.5984\pm0.0035\)（散布 \(0.5935\dots0.6027\)，最小
MC p = 0.36）。GUE 统计对所有结构一致------正如 4.1 节所总结，GUE 的来源
是 AB 云动力学，而非旋量结构的选择。

测试 38 的多语言移植（冻结数据 + 自实现 Jacobi 算法，不用 LAPACK）随包
提供 10 种语言：
\texttt{verification/\{cpp,java,rust,go,fortran,haskell,r,matlab,julia,javascript\}/spinor38/}。

\end{document}
